\section{All indices for the cubic ternary lattice}\label{sec:cubic}
Take $L=\mathbb Z^3$ with the standard dot product. Write
$a(m)$ for the number of intrinsic classes at index $m$ and
$r_3(b)=\#\{(x,y,z)\in\mathbb Z^3:x^2+y^2+z^2=b\}$.
Ordinary and proper classes coincide in rank three. The coefficients
are expressed in terms of $r_3$, quadratic residue symbols and rational
functions of one variable.
For an odd prime $p$ write $\chi_p(b)=(b/p)$. In the following formulas
$T$ is formal; products over $p^k\parallel m$ use only $k>0$.

\paragraph{Local factors at odd primes.}
Put
\[
 F_1=\frac1{(1-T)(1-pT)(1-p^2T)},\quad
 F_2=\frac1{(1-T)^2(1-pT)}.
\]
Also define
\[
 J_p(b)=\frac1{1-T}
 \begin{cases}(1-T)^{-2}&\chi_p(b)=1,\\
 (1-T^2)^{-1}&\chi_p(b)=-1.\end{cases}
\]
The identity row is $F_1(1-(p+1)T/2)$. The half-turn row has an
axis parameter $b$ and a squarefree odd reversal index $e$. At $p\mid e$
it is $TF_2/2$. At $p\nmid e$ it is
$F_2(1-(1+\chi_p(-b))T/2)$. Here $b=e$ or $2e$.
The quarter-turn row is $J_p(-1)(1-(1+\chi_p(-1))T/2)$.
The third-turn and sixth-turn rows at $p\ne3$ are both
$J_p(-3)(1-(1+\chi_p(-3))T/2)$. At $p=3$ they are, respectively,
\[
 1+\frac{T(1+2T)}{2(1-T)^2},\qquad \frac{T}{2(1-T)^2}.
\]
The sixth-turn row is used only when $3\mid m$.

\paragraph{Local factors at two.}
For the identity retain $F_1$ with $p=2$ and set
$P=(1+3T)/((1-T)^2(1-2T))$. The other rows are
\[
\begin{array}{c|c}
 \text{rotation}&\text{factor at }2\\ \hline
 \text{half-turn, }b\text{ odd}&P\\
 \text{half-turn, }b\text{ even}&(1-2T)^{-1}+TP\\
 \text{quarter-turn}&(1+T)/(1-T)^2\\
 \text{third- or sixth-turn}&((1-T)(1-T^2))^{-1}.
\end{array}
\]
For each row $j$ write $W_j(m)=\prod_{p^k\parallel m}[T^k]A_{j,p}(T)$.
Use $W_2(b,e;m)$ for the half-turn row.
\begin{theorem}\label{thm:cubic-all}
For every $m\ge1$,
\[
\begin{split}
24a(m)={}&W_1(m)+6W_4(m)+8W_3(m)
  +8\mathbf1_{3\mid m}W_6(m)\\
 &+\sum_{\substack{e\mid m\\e\ {\rm odd\ squarefree}}}
  \left(\frac{r_3(e)}2W_2(e,e;m)
       +\frac{r_3(2e)}2W_2(2e,e;m)\right).
\end{split}
\]
Here the subscripts $1,4,3,6$ denote rotation orders.
\end{theorem}
In particular $a(2)=a(3)=3$, $a(6)=13$. The cubic Dirichlet series
does not have multiplicative coefficients. For $p\ge5$ the rows give
\[
 a(p)=\frac{2p^2+19p+70+3\chi_p(-1)-6\chi_p(-2)+8\chi_p(-3)}{48}
       +\frac{r_3(p)+r_3(2p)}{96}.
\]
At any fixed prime $p$, the formula gives a rational series in $T$
for the coefficients $a(p^k)$: only $e=1,p$ occur, and at $p=2$
only $e=1$.

\subsection*{Maximal parents and their local trees}
An integral unimodular positive rank-three lattice is isometric to $L$.
The open ball of radius $\sqrt2$ has volume greater than $2^3$,
so Minkowski's theorem gives a norm-one vector. Its orthogonal complement
is integral unimodular; the open disk of radius $\sqrt2$ has area greater
than $2^2$, giving another norm-one vector. Thus every integral self-dual
lattice in the rational cubic space is a possible parent. The positive
integral stabilizer $\Gamma=\mathrm{SO}(L)$ consists of the 24 determinant-one
signed permutations.

Over $\mathbb Q_2$ the norm-integral vectors are exactly
$L_2=\mathbb Z_2^3$: an integral triple with an odd coordinate has sum
of squares nonzero modulo four. Consequently the self-dual parent at two
is unique. For $p$ odd let the vertices be the self-dual lattices in
the given local space; join vertices whenever their intersection has
index $p$ in both. For any two vertices $P,R$ the elementary divisor
exponents relative to one another are $(-h,0,h)$ since duality negates
these exponents. An index-$p$ hyperplane $K\subset P$ corresponds to
a line $\ell=(K/pP)^\perp$ in the nonsingular space $P/pP$. If
$\ell$ is anisotropic the discriminant group $K^*/K$ is cyclic of
order $p^2$: split off a unit vector lifting $\ell$, then $K$ replaces
this summand by its $p$-multiple. There is just one index-$p$ overgroup.
If $\ell$ is isotropic, choose a lift $v\in P$ with $v^2\in p^2\mathbb Z_p$
by adjusting $v$ by $p$ times a vector pairing to a unit with it.
Then $K+\mathbb Z_p(v/p)$ is a second self-dual parent. The two
distinct order-$p$ subgroups of $K^*/K$ defining these parents show
that the group is elementary; its pairing, multiplied by $p$, is a
nondegenerate hyperbolic plane. Its two isotropic lines are all the
possibilities. The neighbors therefore correspond to the isotropic lines. There are $p+1$ lines, since a nondegenerate ternary
form over a finite field of odd order has an isotropic point and the
conic through that point is parametrized by the projective line.
After diagonalization and fixing the third coordinate, isotropy follows
by intersecting two subsets of $\mathbb F_p$, each containing $(p+1)/2$
quadratic values.

Choose a Smith basis of $P$ in which $R$ has generators
$p^h u,v,p^{-h}w$ with $h>0$. The lattice with generators
$pu,v,p^{-1}w$ is integral by the pairings in $P$ and $R$ and has
unit determinant. Continuing with each exponent between 0 and $h$
gives a path of length $h$; its members all contain $P\cap R$.
The exponent $h$ is the least integer such that $p^hR\subset P$,
so it changes by at most one along edges. The line $p^hR\bmod pP$
is spanned by $w$ and is isotropic. Any other isotropic line has
nonzero pairing with $w$, since a nonsingular ternary space contains
no totally isotropic plane. The corresponding neighbor omits $p^hR$,
so its distance exponent to $R$ is $h+1$. Every vertex has a unique edge toward the root, with no edges
within a level; hence the neighbor graph is a tree. For a fixed full lattice $M_p$, the set of parents containing $M_p$
is a finite convex subtree when nonempty. Every vertex on the path
between two such parents contains their intersection, and every parent
lies between $M_p$ and $M_p^*$.

For a represented lattice $M$ of determinant $m^2$ the product of these
subtrees is the complex of its possible parents. Each vertex gives an
actual global parent by intersecting its local lattices with the rational
space. This local-to-global correspondence follows by placing
all lattices between $cL$ and $c^{-1}L$ and decomposing their finite
quotient into its primary groups. The dyadic factor and every $p\nmid m$
contribute one vertex. The index of $M$ in every parent is $m$.

\subsection*{The signed count of marked cubes}
The product of parent trees is contractible. For $g\in\mathrm{SO}(M)$
the alternating sum of its traces on the rational cellular chains is 1:
cancel the traces on boundaries in successive dimensions to leave
only $H_0$. On a preserved $r$-cube the contribution is
$(-1)^r(-1)^{b_g}$ where $b_g$ is the number of reversed local edges.
Prime places cannot be permuted. Averaging over $\mathrm{SO}(M)$
and summing over represented classes, we may group terms by unanchored
cubes. For such a cube $C$, let $K_C$ be the intersection of its
vertices and $S_C$ its rational proper stabilizer. It is finite, since the stabilizer of any one
of its finitely many vertices is finite. If edges occur at the primes
of $d$, its intersection has index $d$ in every vertex. The lattices
with $C$ among their parents are exactly the index-$m/d$ sublattices
of $K_C$. Changing the order of the finite incidence sum, the contribution
of the orbit of $C$ is
\[
 \frac{(-1)^{\omega(d)}}{|S_C|}
 \sum_{g\in S_C}(-1)^{b_g}\,
 \#\{M\subset K_C:[K_C:M]=m/d,\ gM=M\}.
\]
In both orbit decompositions, a marked pair has stabilizer
$S_C\cap\mathrm{SO}(M)$, on which the character $(-1)^{b_g}$ is
evaluated. We may mark one vertex of each cube as $L$.
Summing over all cubes anchored there replaces the preceding denominator
by $2^{\omega(d)}|\Gamma|$. To see this, the sum over $S_C$-orbits
of vertices of $1/|(S_C)_P|$ equals $2^{\omega(d)}/|S_C|$;
after identifying $P$ with $L$ its remaining changes of coordinates
are precisely $\Gamma$, independently of whether the cube admits an
edge-reversing symmetry.

For an anchored cube $gL$ is a corner and $g^2L=L$ because $g$ preserves
or reverses each local edge. Write $e$ for the product of primes where
$gL_p\ne L_p$. At these primes an edge is forced, $K_{C,p}=L_p\cap gL_p$,
and the signs from dimension and orientation cancel. At an odd prime
outside $e$, either there is no edge or one chooses a $g$-fixed isotropic
line at $L_p$. Thus, writing $F_{g,P}(T)$ for the series of $g$-stable
sublattices of $P$ by index exponent, the contribution of $g$ factors
into local coefficients. At a reversed prime the factor is
$TF_{g,L_p\cap gL_p}/2$. At another odd prime it is
\[
 F_{g,L_p}-\frac T2\sum_{\substack{\ell\ \text{isotropic}\\g\ell=\ell}} F_{g,K_\ell}.
\]
At two it is $F_{g,L_2}$. Conversely any rational rotation with square
fixing $L$, with squarefree odd displacement $e\mid m$, and with the
indicated choices at the remaining primes stabilizes exactly that cube.
The product for invariant sublattices follows by primary decomposition
of $K_C/(m/d)K_C$.

\subsection*{Classification of rational rotations}
Write a rational rotation as conjugation on the trace-zero Hamilton
quaternions by a primitive integral quaternion $q=(w,x,y,z)$, unique
up to sign. To see surjectivity, an integral or rational orthogonal
rotation of trace $-1$ is a half-turn about a rational axis, represented
by its pure axis vector. Otherwise its Cayley transform
$(g-I)(g+I)^{-1}$ is rational skew-symmetric and is the cross-product
matrix of a rational vector $v$. The quaternion $(1,v)$ gives back
$g$ by multiplying out the Cayley transform. Central scalars are exactly
the kernel. The least denominator of the rotation matrix is the odd part
of $q\bar q$. Indeed its entries have this norm as common denominator,
with diagonal numerators $w^2+x^2-y^2-z^2$ and their cyclic translates
and off-diagonal numerators $2(xy\pm wz)$ and their cyclic translates.
At an odd prime dividing the norm these cannot all vanish: vanishing
forces the six pairwise products of the four coordinates to vanish
modulo $p$, and the zero norm then forces every coordinate to vanish.
At two the primitive norm has valuation at most two. If it has valuation
one all numerators are even; if it has valuation two all four coordinates
are odd, and all numerators are divisible by four. Because of the
symmetric three Smith exponents the least denominator is also
$[L:L\cap gL]$. This agrees with the cubic coincidence index \cite{Baake}.

The denominator here is squarefree and equals $e$.
As $q^2=2wq-q\bar q$, its coefficients are divisible by $p\mid e$
if and only if $p\mid w$. The square rotation is integral precisely
when this holds for every such $p$: divide $q^2$ by $p$ to remove the
entire $p$-part of its norm. For $e\ge5$ the estimate
$|w|\le2\sqrt e<e$ forces $w=0$. Pure primitive quaternions have
norm $e$ or $2e$, since a sum of three squares divisible by four
has all coordinates even. Conversely every triple of norm $e$ or
$2e$ is primitive by squarefreeness. These give $r_3(e)/2,r_3(2e)/2$
rotations. For $e=3$ there are also the eight rotations represented
by $(3,\pm1,\pm1,\pm1)$, of order six. For $e=1$ there are the 24
rotations of the cube: the identity, nine half-turns (already included
by the pure quaternions), six quarter-turns, and eight third-turns.

\subsection*{Evaluation of the local traces}
At an odd prime the identity fixes all $p+1$ conic points. A half-turn
with primitive axis of norm $b$ at $p\nmid 2b$ has an anisotropic
fixed axis; the determinant of the orthogonal plane has square class
$b$. Hence $1+\chi_p(-b)$ isotropic lines are fixed. For a quarter-turn,
or a third- or sixth-turn at $p>3$, the axis is again anisotropic.
The other two eigenlines are isotropic whenever defined since their
eigenvalues square to a number different from one. The counts are
$1+\chi_p(-1)$ and $1+\chi_p(-3)$ respectively. The quarter-turn calculation also applies at $p=3$. An integral third-turn is conjugate by a signed permutation
to the cyclic permutation of the three coordinates. It fixes the unique eigenline $(1,1,1)$ modulo three,
which is isotropic.

At primes not dividing the rotation order the action decomposes by
the distinct irreducible residue factors of the minimal polynomial:
the Bezout projections lift over $\mathbb Z_p$. For a half-turn the
summands have ranks one and two, and the unrestricted subgroup series
on these summands give $F_2$. For the other rotations with two non-real
eigenvalues, the nontrivial plane either splits into two lines or is
free of rank one over the unramified quadratic integer ring. Its stable
sublattices then give $(1-T)^{-2}$ or $(1-T^2)^{-1}$. This description
applies equally to each stable edge intersection. For the identity use the
rank-three HNF series $F_1$. For the exceptional order-six rotations at
three the polynomial $X-1$ is coprime to $X^2-X+1$ over $\mathbb Z_3$.
The plane is rank one over $\mathbb Z_3[\zeta_6]$, a discrete valuation
ring with residue field of size three: $\zeta_6+1$ is a uniformizer
by its Eisenstein polynomial at three. Thus its factor, including the
axis, is $(1-T)^{-2}$, also on the reversed edge.

At the exceptional primes, the local module calculation reduces
to a fiber-product order. Let $O$ be a local integer ring with
residue field $\mathbb F_p$ and uniformizer $\pi$. Consider the order
\[
 A=\{(x,y)\in\mathbb Z_p\times O:x\bmod p=y\bmod\pi\}.
\]
For an $A$-stable lattice $E$ with saturated module
$(\mathbb Z_p\times O)E=D=p^a\mathbb Z_p\oplus\pi^bO$, one has
$(p,\pi)D\subset E$ (here $(p,\pi)D$ denotes the direct sum of the
two maximal-ideal multiples). The image of $E$ modulo this direct sum
projects onto both residue spaces. The possibilities are the whole
plane or one of the $p-1$ graphs of nonzero maps. Conversely all these
possibilities work. For ambient $\mathbb Z_p\oplus O$ the series is
$F_{\rm split}=(1+(p-1)T)/(1-T)^2$. For ambient $A$, when $a,b\ge1$,
all graphs and the plane occur, with one power of $T$ removed because
the ambient has index $p$. When only one of $a,b$ is zero, the
residue-matching condition is impossible; at $(0,0)$ only the graph
defining $A$ occurs. Thus $F_{\rm glued}=1+TF_{\rm split}$.
For the third-turn at three, evaluation at $(1,\zeta_3)$ identifies
the permutation lattice $\mathbb Z_3[X]/(X^3-1)$ with $A$ for
$O=\mathbb Z_3[\zeta_3]$ and $\pi=1-\zeta_3$.
The two factors have resultant $3$ and common residue field
$\mathbb F_3$; polynomial representatives therefore give pairs
with the same residue. The hyperplane for the fixed isotropic line imposes
coordinate sum zero modulo three, so its module is $3\mathbb Z_3
\oplus\pi O$, isomorphic to the split module. This gives
$F_{\rm glued}-TF_{\rm split}/2$. At two the quarter-turn is
the direct sum of its fixed axis and $\mathbb Z_2[i]$; evaluation
at $(1,i)$ gives the same order $A$ with $\pi=1-i$ owing to the
resultant $2$, and the ambient is split. The remaining third- and
sixth-turns at two have a separate axis and the unramified quadratic
plane.

Finally, a dyadic half-turn has saturation
$R=\mathbb Z_2\oplus\mathbb Z_2^2$ under the two rational eigenprojections.
In coordinates on its projected summands the original lattice is $R$
or the graph of a nonzero linear form on the negative residue plane:
it contains $2R$, the two projections are surjective, and the only
other subspace with surjective projections is the full residue space.
The split case is equivalent to integral eigenprojections, or
$g\equiv I\pmod2$. For a primitive axis $v$ with norm $b$ the half-turn
is $2vv^t/b-I$, so this is exactly the case of odd $b$.
A sublattice with saturation $2^a\mathbb Z_2\oplus D$ for
$D\subset\mathbb Z_2^2$ has four choices (one full residue space and
three graphs), giving the series $P$ for a split ambient.
In a glued ambient, when $a\ge1$ one must have
$D\subset\ker(\ell\bmod2)$ for its defining form $\ell$, and then
every choice works, contributing $TP$. When $a=0$, $D$ must avoid
containment in this kernel; exactly one graph works. The series of
such $D$ is $(1-T)/((1-T)(1-2T))$. The asserted dyadic factors follow.

\subsection*{A branch point inside the domain of continuation}
Let $Z_{\rm cub}(s)=\sum_m a(m)m^{-s}$ and let all powers on the real
line initially have their positive values. Define
\[
 R(s)=(1-2^{1-s})^{-1/2}
 \prod_{p>2}\big(1-\tfrac12(p+1)p^{-s}\big)
                  (1-p^{1-s})^{-1/2}.
\]
For $\Re s>3/2$ the product is nonvanishing and holomorphic:
on compact subsets the logarithm of the $p$-factor is
$O(p^{-\Re s}+p^{2-2\Re s})$, and each factor is nonzero.
For a positive integer $b$ put $P_b(X)=1+3X$ if $b$ is odd and
$P_b(X)=1-X+4X^2$ if $b$ is even. Set
\[
 B(s)=\sum_{e\ {\rm odd\ squarefree}}
 \frac{e^{-s}}{2^{\omega(e)}}\sum_{b\in\{e,2e\}}r_3(b)P_b(2^{-s})
  \prod_{p\nmid 2e}\left(1-\frac{1+\chi_p(-b)}2p^{-s}\right).
\]
This series is normally convergent on $\Re s>3/2$. The products over
$p$ are bounded uniformly in $b,e$ on every compact subset of
$\Re s>1$, and $\sum_{n\ge1}r_3(n)n^{-\sigma}$ converges for
$\sigma>3/2$ by summing points of $\mathbb Z^3$ in dyadic annuli.
One has $B(2)>0$.

\begin{corollary}\label{cor:cubic-boundary}
There is a holomorphic function $H$ on $\Re s>1$ such that
\[
 Z_{\rm cub}(s)=
 \frac{\zeta(s)\zeta(s-2)\zeta(s-1)^{1/2}R(s)}{24}
 +\frac{\zeta(s)^2\zeta(s-1)B(s)}{48}+H(s).
\]
The identity holds initially for $\Re s>3$. It gives a meromorphic
continuation to $\Re s>2$ and continuation on paths in
$\Re s>3/2$ avoiding zeros and the pole of $\zeta(s-1)$, with the
usual continuation past isolated meromorphic poles. The pole at three
is simple, with residue
\[
 \rho=\frac{\zeta(3)\zeta(2)}{24}
     \prod_{p>2}\left(1-\frac{p+1}{2p^3}\right).
\]
At two, in a slit disk reached from the right,
\[
 Z_{\rm cub}(s)=\frac{\zeta(2)^2 B(2)}{48(s-2)}
       -\frac{\zeta(2)R(2)}{48\sqrt{s-2}}+O(1).
\]
In particular no meromorphic germ at two agrees with this continuation.
\end{corollary}
\begin{proof}
For the identity row, factor $\zeta(s)\zeta(s-1)\zeta(s-2)$
out of the subgroup series; its odd-prime correction is
$\prod_{p>2}(1-(p+1)p^{-s}/2)$.
Extracting $\zeta(s-1)^{-1/2}$ leaves $R$.
For every half-turn row put back the three common denominator factors
of $\zeta(s)^2\zeta(s-1)$, including at the flipped primes. The remaining
Euler factors at $p\mid e$ produce exactly $e^{-s}2^{-\omega(e)}$;
the dyadic factors give $P_b$. The quarter-, third- and sixth-turn
products are holomorphic on $\Re s>1$ by absolute convergence, since
their good factors differ from one by $O(p^{-\Re s})$ uniformly on
each compact subset. They give $H$ (including their group-order
weights). At two use $\zeta(0)=-1/2$ and
$(s-2)\zeta(s-1)=1+O(s-2)$.
\end{proof}

The decomposition also gives the unsmoothed estimate
\[
 \sum_{m\le X}a(m)=\frac{\rho}{3}X^3+
 O\big(X^2(\log(2X))^{3/2}\big).
\]
For the error term, the summatory coefficients of
$\zeta(s)\zeta(s-1)\zeta(s-2)$ are
$\zeta(3)\zeta(2)X^3/3+O(X^2\log(2X))$: sum the squares of positive
integers up to $X/(a b)$ with weight $b$ and bound the missing terms
of $\sum_{a,b}a^{-3}b^{-2}$ by $O(1/X)$.
The unsigned coefficients $h$ of the odd-prime correction satisfy
\[
 \sum_{j\le X}\frac{|h(j)|}{j^2}\le
 \prod_{2<p\le X}\big(1+(p+1)/(2p^2)\big)\ll\sqrt{\log(2X)}
\]
by Mertens' prime reciprocal estimate. Dyadic decomposition gives
\[
 \sum_{j>X}|h(j)|/j^3\ll X^{-1}\sqrt{\log(2X)}.
\]
These estimates
give the claimed error for the identity row. For the half-turn part
strip off $\zeta(s-1)$; the remaining Dirichlet expansion is absolutely
summable at two, by the same absolute bounds used for $B$ and for
$\zeta(2)^2$. Since $\sum_{j\le y}j\ll y^2$, its contribution is
$O(X^2)$. Absolute convergence at two gives the same bound for $H$.

\subsection*{An intrinsic class-number interpretation}
For any $m$, $a(m)$ is also the number of integral positive ternary
Gram classes of determinant $m^2$ with Hasse invariant $+1$ at every
odd $p\mid m$. Integrality here is with respect to the bilinear Gram matrix.
At odd $p\nmid m$, diagonalization of a unimodular matrix gives
unit entries and Hasse invariant $+1$. The real invariant is $+1$,
so Hilbert reciprocity gives the invariant $+1$ at two as well.
Hasse--Minkowski and the square determinant put the lattice in the
rational cubic space. Every full integral lattice in that space
enlarges to a self-dual lattice. At two use norm-integrality as above.
At odd $p$, a maximal integral enlargement exists inside the dual of
the original lattice. On diagonalizing this maximal lattice no exponent
can be $\ge2$, since then division of the corresponding basis vector by
$p$ gives an integral enlargement. The number of exponent-one entries
is even. If it is two the diagonal is $u,pv,pw$ with $u,v,w$ units;
its Hasse invariant is $\chi_p(-vw)$ by the odd-prime Hilbert symbol
formula. Thus the last two unit coefficients define a split residual
plane. An isotropic primitive combination of its two vectors has
norm in $p^2\mathbb Z_p$ and pairs into $p\mathbb Z_p$ with the
lattice, again permitting enlargement. Thus every maximal lattice
at odd $p$ is unimodular. Gluing gives the required unimodular
global parent and unimodular uniqueness identifies it with $L$.
The resulting index is $m$ by the determinant formula. All implications
reverse for a represented lattice.

\subsection*{Other one-class ternary genera}
The tree calculation extends to a ternary parent whose proper
genus has class number one. Let $G$ be a positive integral symmetric
ternary Gram matrix with $D=\det G$, $L$ its lattice, and
$\Gamma=\mathrm{SO}(L)$. Fix a finite set $S$ of primes containing all
divisors of $6D$. Suppose that every lattice in the proper genus of $L$
is properly isometric to $L$. Write $Z_L^{(S)}(s)$ for the subform
Dirichlet series restricted to indices prime to $S$ and
$\zeta^{(S)}(s)=\zeta(s)\prod_{p\in S}(1-p^{-s})$.
\begin{corollary}\label{cor:ternary-genus-one}
For $(m,\prod_{p\in S}p)=1$ the coefficient is the following finite sum.
Range over rational rotations $g$ with $g^2\in\Gamma$, whose least
denominator in the chosen basis is squarefree and divides $m$; call this
denominator $e$. Add $|\Gamma|^{-1}$ times the product over
$p^k\parallel m$ of the $T^k$ coefficients of
\[
 \begin{cases}
 TF_g/2 &p\mid e,\\
 F_g(1-i_p(g)T/2)&p\nmid e.
 \end{cases}
\]
Here $i_p(g)$ counts the $g$-fixed isotropic lines of $L/pL$.
The common lattice-independent stable factor $F_g$ is $F_1,F_2,
J_p(-1),J_p(-3)$ according as the rotation has order 1, 2, 4, or 3
or 6. In the second row $i_p(1)=p+1$; for a half-turn of primitive
axis $v$,
$i_p(g)=1+\chi_p(-D\,v^tGv)$, and for the remaining orders
$i_p(g)=1+\chi_p(-1)$ or $1+\chi_p(-3)$ respectively.
The series is meromorphic on $\Re s>2$ and has a simple pole at three
with positive residue
\[
 |\Gamma|^{-1}\prod_{p\in S}(1-p^{-1})\,
 \zeta^{(S)}(3)\zeta^{(S)}(2)
 \prod_{p\notin S}\big(1-(p+1)/(2p^3)\big).
\]
It has no meromorphic germ across the real point $s=2$.
\end{corollary}
\begin{proof}
At $p\notin S$ all vertices of the local self-dual tree are locally
isometric to $L_p$. Every nonsingular
ternary form over $\mathbb F_p$ contains a hyperbolic plane, and its
remaining line is determined up to squares by the determinant. A
residue-field isometry lifts to $\mathbb Z_p$: given an isometry
modulo $p^j$, a correction by $I+p^jX$ solves the next linearized
Gram equation because $X^tA+AX$ maps onto the symmetric matrices
for an invertible $A$ and odd $p$. A change of determinant sign can
be made by $-I$. Keep the parent fixed at all $p\in S$. Each
resulting global vertex is in the proper genus of $L$ and hence
is a permissible parent. For each $M$ the local tree Euler count
therefore still counts just one isometry class. The local Smith
argument identifies displacement with denominator at $p\notin S$;
at all other primes there is no displacement.

Such a $g$ has finite order because $g^2\in\Gamma$. Rationality of
the characteristic polynomial implies order $1,2,3,4$ or $6$.
The local calculations are tame as $p\nmid6$, and give the displayed
stable factors. For the half-turn the orthogonal plane has determinant
class $D/(v^tGv)$; the resulting symbol gives the stated number of
lines. For the other nonidentity orders the fixed axis is nondegenerate
and the two nontrivial eigenlines are isotropic exactly when present.
All rotations with $g^2\ne1$ form a finite list: a positive real
rotation with a prescribed nonidentity square is supported on the
same axis and has at most two possible angles. Half-turns are enumerated
by primitive integral vectors $v$ up to sign, through
$g=2v(v^tG)/(v^tGv)-I$. If $q=v^tGv$, then
$q\le 2De$: the gcd of the entries of $v(v^tG)$ equals the gcd
of the entries of $v^tG$, and multiplying by the adjugate of $G$
shows this gcd divides $D$. This gives a finite search bound for each
index. At primes dividing $e$, $g^2L=L$ and the endpoints do form
the required edge by the squarefree Smith exponents.

For analytic continuation, the finite set of rotations with $g^2\ne1$
has Euler products absolutely convergent on $\Re s>1$. For each
half-turn strip off $\zeta^{(S)}(s)^2\zeta^{(S)}(s-1)$. The remaining
term is
\[
 2^{-\omega(e)}e^{-s}
 \prod_{p\notin S,\ p\nmid e}(1-i_p(g)p^{-s}/2).
\]
Only squarefree $e$ outside $S$ occur. This series of terms is
normally convergent on $\Re s>3/2$ by $q\le2De$ and the ternary
lattice point bound. Thus all nonidentity contributions are meromorphic
in a neighborhood of two. The identity term is
\[
 |\Gamma|^{-1}\zeta^{(S)}(s)\zeta^{(S)}(s-1)\zeta^{(S)}(s-2)
       \prod_{p\notin S}(1-(p+1)p^{-s}/2).
\]
As above its product is $\zeta^{(S)}(s-1)^{-1/2}$ times a holomorphic
unit near two. Since $\zeta^{(S)}(s-2)$ has zero of order $|S|$ at
two and $\zeta^{(S)}(s-1)$ has a simple pole there, the identity
contribution is $(s-2)^{|S|-1/2}$ times a holomorphic unit, which
precludes a meromorphic germ. Evaluating at three gives the residue
as stated.
\end{proof}
\subsection*{All indices for the lattice of determinant two}
Write $Z_{112}$ for the unrestricted subform series of the lattice
with diagonal Gram matrix $(1,1,2)$. The dyadic parent change gives the following local factors.

\begin{theorem}\label{thm:det-two}
Use Theorem~\ref{thm:cubic-all} with its odd-prime rows and rotation
multiplicities unchanged. At two, include the coefficient of
$T^{v_2(m)}$ for every $m$, including odd $m$, using this table with
$F=((1-T)^2(1-2T))^{-1}$:
\[
\begin{array}{c|l}
\text{rotation}&\text{factor at two}\\ \hline
1 &(3-2T)F_1\\
2,\ b\ {\rm even}&(1+3T)F\\
2,\ b\equiv1\pmod4 &(3-T+2T^2)F\\
2,\ b\equiv3\pmod4 &(3-2T)(1+3T)F\\
4 &(1+T)/(1-T)^2\\
3,\ 6 &T/((1-T)(1-T^2)).
\end{array}
\]
The result is the coefficient of $Z_{112}$ for every $m\ge1$.
Its residue at $s=3$ is $11/4$ times the cubic residue; the same
unsmoothed error bound holds. The coefficient of $(s-2)^{-1/2}$ in
its continuation is $5/2$ times the corresponding cubic coefficient,
and the term of order $(s-2)^{-1}$ has strictly positive coefficient.
\end{theorem}
\begin{proof}
In the original lattice the vectors $(1,1,0),(1,-1,0),(0,0,1)$
give an index-two lattice $K$ of Gram matrix $2I_3$. Use its basis
and write $B=2\,\cdot$, twice the usual dot product. The local lattice
$K_2$ is intrinsic, being exactly the vectors with $B(x,x)\in2\mathbb Z_2$.
Every integral bilinear dyadic lattice $P$ of determinant valuation
one lies in $\frac12K_2$: a primitive cubic vector has norm valuation
at most one. The nonzero residues $2P\bmod 2K_2$ have weight two
since $B(x,x)$ is integral. Distinct weight-two residues have odd dot
product, so their halves cannot both belong to $P$. Covolume then
forces $K_2\subset P$ and $[P:K_2]=2$. Conversely each of the three
weight-two lines $\lambda$ gives an integral parent
$P_\lambda=K_2+\mathbb Z_2\lambda/2$, where any integral lift can be used.
All have pairwise intersection $K_2$.

Every positive integral lattice of determinant two is isometric to
the stated parent. The open ball of radius $\sqrt2$ has volume greater
than $8\sqrt2$, so Minkowski gives a unit vector. Its rank-two
orthogonal complement still has determinant two, and the disk of the
same radius has area greater than $4\sqrt2$. Splitting off another
unit gives the assertion. At odd $p$ the permissible local parents
are exactly the self-dual vertices for $B/2$. Consequently an abstract
represented lattice $M$ in this rational space can be counted by the
same contractible product of odd-prime parent trees as in the cubic
proof, with one dyadic membership test: $M_2$ must lie in at least
one $P_\lambda$. Indeed any permissible choices of odd vertices and
one such dyadic parent glue to an integral lattice of determinant
two. Conversely any embedding in the stated parent supplies these
local choices. To anchor a corner of the odd complex, glue its odd
vertices to $K_2$; the resulting lattice for $B/2$ is integral
unimodular at every prime, hence cubic. This gives exactly the same
anchoring group, rotations and odd local factors as before.

For a dyadic rotation $g$ write $F_{g,E}$ for the stable-sublattice
series in $E$. For $M_2\subset K_2$, its index in any $P_\lambda$
is twice its index in $K_2$. Every other permissible $M_2$ belongs
to just one $P_\lambda$, which must be fixed if $M_2$ is fixed.
Thus the requisite trace is
\[
 T F_{g,K_2}+
 \sum_{\lambda:\,gP_\lambda=P_\lambda}
       (F_{g,P_\lambda}-T F_{g,K_2}).
\]
This counts the identity once for every eligible lattice at this place,
regardless of how many parents contain it.

The identity fixes all three parents. A quarter-turn permutes two
coordinates modulo two and fixes just the line they span by their sum.
This enlarges its Gaussian plane by the fractional ideal
$(1-i)^{-1}\mathbb Z_2[i]$, leaving the module type unchanged.
The other rotations of orders three and six have an unramified
irreducible two-dimensional summand modulo two, whose nonzero elements
are precisely the three weight-two vectors: the fixed axis is the
nonisotropic line and no weight-two line is fixed. Indeed the weight-two
plane is the kernel of the invariant coordinate-sum form and is the
unique nontrivial residue summand.

For a half-turn of primitive axis $v$ let $E=vv^t/b$ be the positive-axis
projection. A stable ambient is split if and only if $E$ preserves it;
its two possible series are $S=(1+3T)F$ and $G=(1-2T)^{-1}+TS$
from the cubic calculation. If $b$ is even then the reduction of $g$
interchanges the two coordinates on which $v$ is odd. Its only fixed
parent adjoins $v/2$ to $K_2$. The image $E K_2$ lies in
$\mathbb Z_2v/2$, so the enlarged parent is split. For odd $b$,
$K_2$ is split, all three parents are fixed and $P_\lambda$ is split
exactly when $E\bmod2$ preserves $\lambda$. When $b\equiv1\bmod4$,
$v$ has weight one modulo two: precisely the line avoiding that
coordinate works, giving $S+2G-2TS$. When $b\equiv3\bmod4$, $v$
has weight three and is orthogonal to each weight-two residue, giving
$(3-2T)S$. These are the rows in the table.

The analytic calculation is unchanged except for these explicit
dyadic factors. The identity term acquires $3-2^{1-s}$; at $s=3$
and $s=2$ this equals $11/4$ and $5/2$. The half-turn polynomial
after removing $F$ belongs to the displayed finite list, and is
positive at $T=1/4$ in all cases. The same absolute estimates therefore
hold for its meromorphic contribution, with positive residue at two.
The finitely many remaining rotations are holomorphic near two, and
the finite dyadic change preserves the summatory error bound.
\end{proof}
For the determinant-two parent, the coefficients at indices
$5,7,35$ are respectively $9,13,234$.
\subsection*{A finite test for dyadic changes of the cubic parent}
Let $C=\mathbb Z^3$ with the cubic metric, and take any lattice $P$
with $2^aC\subset P\subset C$. A common scalar change of the metric or lattice leaves the
construction unchanged. Let $\mathcal P$
be the orbit of $P_2$ under $\mathrm{SO}(C_2)$. Suppose that the
integral signed-permutation rotation group $\Gamma$ acts transitively
on $\mathcal P$. The lattices are
given by subgroups of $C/2^aC$, and the needed residue matrices of
$\mathrm{SO}(C_2)$ can be listed by primitive quaternion coordinates
modulo $2^{a+2}$. In fact every dyadic rational rotation is given by
a quaternion by the Cayley argument above over $\mathbb Q_2$; a
primitive quaternion norm has valuation at most two. Dividing the
matrix numerators and the norm by the common power of two therefore
determines the rotation modulo $2^a$. Each primitive residue lift is
a genuine dyadic rotation. For example, modulo two the group consists
exactly of the coordinate permutations: a vector of norm one has
weight one or three, and the weight-three vector cannot be a column
of an orthonormal basis as its complement is alternating. The integral
rotations realize all permutations. The orbit hypothesis consequently
holds for \emph{every} intermediate $2C\subset P\subset C$.

\begin{corollary}\label{cor:dyadic-cubic-family}
Under the orbit hypothesis the full subform series has a finite
coefficient formula at every index, with explicit leading residue,
unsmoothed error term and real square-root branch at two.
Keep the odd local factors and the individual rational rotations in
the cubic formula, grouping half-turns only when all their local
coefficients coincide. For each rotation $g$ replace its dyadic row
by the following factor even when the index is odd:
\[
 D_g(T)=
 \sum_{\substack{\varnothing\ne J\subset\mathcal P\\gJ=J}}
 (-1)^{c_g(J)+1}\,
 T^{\log_2[P_0:E_J]}F_{g,E_J}(T).
\]
Here $E_J=\bigcap_{P'\in J}P'$, $P_0$ is any member of $J$, and
$c_g(J)$ counts the cycles of the permutation induced by $g$ on $J$.
Every factor $F_{g,E_J}$ is one of the dyadic series evaluated in
the preceding proofs. The original cubic residue at three and the
coefficient of $(s-2)^{-1/2}$ are multiplied, respectively, by
$A(1/8)$ and $A(1/4)$, where $A(T)=D_1(T)/F_1(T)$ is a polynomial
positive at both arguments. The counting error remains
$O(X^2(\log(2X))^{3/2})$.
\end{corollary}

\begin{proof}
Glue any collection of odd-prime self-dual parents to the intrinsic
dyadic cubic lattice $C_2$. This gives a global positive unimodular
lattice, necessarily cubic. In its coordinates any further choice
$P'\in\mathcal P$ corresponds to the stated parent $P$ by the assumed
transitivity of the integral rotation group. Hence the odd-prime
parent complex of a represented $M$ is unchanged, while dyadic
membership is precisely $M_2\subset P'$ for some $P'\in\mathcal P$.

If $gM_2=M_2$ the permutation $g$ preserves its set of permissible
dyadic parents. On any nonempty finite set with $t$ such cycles, the
alternating sum over nonempty unions of cycles is
$\sum_{j=1}^t(-1)^{j+1}\binom tj=1$; for an empty set it is zero.
Interchanging the resulting finite sums gives the displayed polynomial
combination of stable series. For half-turns the series is the split
or glued half-turn series according as the positive eigenprojection
preserves $E_J$ or not. For a quarter-turn project to the axis and to
the Gaussian plane. The two projected modules have ranks one over
$\mathbb Z_2$ and $\mathbb Z_2[i]$; the fiber-product classification
above gives a split ambient or a glued ambient according as the axis
projection preserves $E_J$ or not. For orders three and six there is
always a split axis and an unramified quadratic plane. Eigenprojections
can be tested on the generators of
$E_J$ by rational arithmetic and dyadic valuations.

Each half-turn row after removing $F_2$ has a polynomial drawn from
a finite list, since $\mathcal P$ is fixed. At $T=1/4$ it is positive:
the genuine local stable-lattice series converges there and contains,
for example, a sufficiently small scalar multiple of $C_2$ in every
parent. The global half-turn sum satisfies the preceding normal-convergence
estimate before axes are grouped by norm. The identity
row is $A(T)F_1(T)$. Its coefficients dominate those of $F_1(T)$
because any one of the parents by itself already contributes that
many sublattices. Hence $A(T)\ge1$ on $0<T<1/4$, and $A(1/4)\ge1$
by continuity. The singular expansion and the error estimate now follow
by the proof of Corollary~\ref{cor:cubic-boundary}.
\end{proof}
\begin{corollary}\label{cor:fcc-bcc}
The face-centered cubic and body-centered cubic lattices have identical
proper and ordinary subform Dirichlet series. They are not similar.
\end{corollary}
\begin{proof}
Use the face-centered lattice $D=\{x\in C:\sum x_i\equiv0\bmod2\}$
and the integral representative of its dual
$E=2D^*=\{x\in C:x_1\equiv x_2\equiv x_3\bmod2\}$ for the
body-centered lattice. The dual description follows because $2C\subset
D$ and its residue plane has orthogonal complement spanned by $(1,1,1)$.
Each has only one dyadic parent in its rotation orbit since every
residue rotation permutes the coordinates. At odd places their local
lattices both equal $C_p$. For each half- or quarter-turn the dyadic series
depends just on whether its rational axis projection preserves the
ambient lattice. This projection $U$ is self-adjoint for the dot
product, so $UD_2\subset D_2$ if and only if $UD_2^*\subset D_2^*$,
by the definition of the dual and reflexivity of full lattices.
Scaling by two does not change this property. The dyadic series of the
other rotation orders are independent of the stable ambient. Thus
each local factor agrees in the two computations of
Corollary~\ref{cor:dyadic-cubic-family}. There are twelve shortest
vectors in $D$ of squared length two and eight in $E$ of squared
length three; a similarity preserves the number of shortest vectors.
\end{proof}
